%%%PAGE 291 rot=0 legibility=good summary=剪贴的两道印刷题及手写解答：2018全国III卷25题（带电小球抛出，重力场+匀强电场中的动能与电势能变化）；水平面上受恒定拉力并带电的小球通过A、P、B三点（余弦定理、电势差）
%%%BLOCK id=291-01 ch=09 sec=电场中的功能关系·重力场与匀强电场叠加 kind=problem alt=04,06 cont=none y=0.02-0.47
% 题干为剪贴的印刷题，题干上有少量手画的括号/斜线标记；题图旁（剪贴纸右下角）有手写批注
\begin{problem}
（20分）如图，$O$、$A$、$B$ 为同一竖直平面内的三个点，$OB$ 沿竖直方向，$\angle BOA=60^\circ$，$OB=\dfrac32OA$。将一质量为 $m$ 的小球以一定的初动能自 $O$ 点水平向右抛出，小球在运动过程中恰好通过 $A$ 点。使此小球带电，电荷量为 $q$（$q>0$），同时加一匀强电场，场强方向与 $\triangle OAB$ 所在平面平行。现从 $O$ 点以同样的初动能沿某一方向抛出此带电小球，该小球通过了 $A$ 点，到达 $A$ 点时的动能是初动能的 3 倍；若该小球从 $O$ 点以同样的初动能沿另一方向抛出，恰好通过 $B$ 点，且到达 $B$ 点时的动能为初动能的 6 倍。重力加速度大小为 $g$。求

(1) 无电场时，小球到达 $A$ 点时的动能与初动能的比值；

(2) 电场强度的大小和方向。

%FIG p291_a 0.478 0.062 0.58 0.152
\notefig[0.25]{figures/p291_a.png}

手写批注：(1) $7:3$\quad (2) \unclear{…}% CHECK: 剪贴纸右下角手写“(1) 7:3”，其下“(2)”之后无法辨认（笔迹与印刷行“的比值；”相交）
\begin{solution}
(1)
\begin{align*}
d\sin60^\circ&=V_0t\\
d\cos60^\circ&=\tfrac12gt^2\\
E_{K0}&=\tfrac12mV_0^2\\
E_{K0}&=\tfrac38mgd\\
E_{KA}&=E_{K0}+\tfrac12mgd\\
\frac{E_{KA}}{E_{K0}}&=\frac73
\end{align*}

(2)
\begin{align*}
\Delta E_{pA}&=3E_{K0}-E_{K0}-\tfrac12mgd=\tfrac23E_{k0}\\
\Delta E_{pB}&=6E_{K0}-E_{0}-\tfrac32mgd=E_{k0}
\end{align*}
% CHECK: 第二式中第二个 E 的下标只写了 0（应指 E_{K0}）

因匀强电场中电势\unclear{均}匀降落

设 $M$ 与 $A$ 等电势
\begin{align*}
\frac{MO}{\frac32d}&=\frac{\Delta E_{pA}}{\Delta E_{pB}}\\
\therefore MO&=d\qquad \alpha=30^\circ\\
qEd\cos\alpha&=\Delta E_{pA}\\
\therefore E&=\frac{\sqrt{3}\,mg}{6q}
\end{align*}
% 末行分母下半被下方剪贴纸遮挡，只露出上半截（读作 6q）；“∴ MO=d”之后的“α”与“d”手写相近，按图中角标读作 α
\end{solution}
\end{problem}
%%%END
%%%BLOCK id=291-02 ch=09 sec=电场中的功能关系·恒力与匀强电场叠加 kind=problem alt=04,06 cont=none y=0.46-0.97
% 剪贴的印刷题；原稿手工框出：“光滑绝缘”“PA=4PB”“质量为 m”；题干中另有手写的圈号①②③（约在“倍”字上方、“沿某一方向抛出”和“沿另一方向抛出”附近），以及“倍。”后一条向下的大弧线；题图上有大量手绘标注
\begin{problem}
（20分）如图，在\fbox{光滑绝缘}水平面上有 $A$、$P$、$B$ 三个点，\fbox{$PA=4PB$}。一\fbox{质量为 $m$}的小球受到沿 $PA$ 方向的恒定拉力作用，并从 $P$ 点在水平面内以某一初速度垂直 $PA$ 抛出，小球在运动过程中恰好通过 $B$ 点，此时其速度大小为初速度大小的 $\dfrac{\sqrt{13}}{2}$ 倍①。使小球带上电荷量为 $+q$ 的正电，同时加上一水平方向的匀强电场，若将此带电小球从 $P$ 点以同样的初速度大小在水平面内沿某一方向抛出②，小球也通过 $B$ 点，此时小球的速度大小是初速度大小的 2 倍，若该带电小球从 $P$ 点以同样的初速度大小在水平面内沿另一方向抛出③，恰好通过 $A$ 点，此时小球的速度大小为初速度大小的 $\sqrt{17}$ 倍。已知小球无论是否带电，均受到大小为 $F$、沿 $PA$ 方向的恒定拉力作用。求：

(1) $\angle APB$ 的正切值；

(2) 电场的场强大小以及场强方向与 $PA$ 的夹角的余弦值；

(3) 有电场且小球带电时，小球从 $P$ 点运动到 $A$ 点过程中与从 $P$ 点运动到 $B$ 点过程中的电势能变化量的比值。

%FIG p291_b 0.34 0.598 0.585 0.742
\notefig[0.3]{figures/p291_b.png}
\begin{solution}
(1)\quad $\left\{\begin{aligned} L\cos\theta&=\tfrac12at^2\\ L\sin\theta&=V_0t\\ V^2&=V_0^2+(at)^2\\ \tfrac{V}{V_0}&=\tfrac{\sqrt{13}}{2}\end{aligned}\right.$\qquad $\tan\theta=\dfrac43$

(2)\quad 不加电场时\ $FL\cos\theta=\tfrac12mV^2-\tfrac12mV_0^2$\ $\to$\ $FL=\dfrac{15}{8}mV_0^2$

加电场\ $\left\{\begin{aligned} U_{PB}q+FL\cos\theta&=\tfrac12mV_B^2-\tfrac12mV_0^2\\ U_{PA}q+F4L&=\tfrac12mV_A^2-\tfrac12mV_0^2\\ V_B=2V_0&\quad V_A=\sqrt{17}\,V_0\end{aligned}\right.$

$\to U_{PB}=\dfrac{3mV_0^2}{8q}$\qquad $U_{PA}=\dfrac{mV_0^2}{2q}$
% “FL=15/8 mV_0^2”右端有一条长弧线箭头一直指向“→U_{PB}=…”那一行（表示把 FL 代入）

由余弦定理\ $BC=4\sqrt{\dfrac25}\,L$

\unclear{由正}% 末行仅两个字，字形不清（可能是“由正（弦定理）”未写完）
\end{solution}
\end{problem}
%%%END
%%%PAGEEND 291
